\chapter{Introdução}

A área da Computação Gráfica tem como objetivo, em geral, a criação algorítmica
de imagens.

Nesse trabalho, nos preocupamos com o problema de síntese de imagens
realísticas, em que, baseando-nos nas leis da física, tentamos criar uma imagem
semelhante à que observamos em fotos e câmeras da vida real. Ao longo do tempo,
desde sua criação, diversos algoritmos foram propostos, cada um com níveis de
abstração diferentes.

O \emph{Path Tracing} é uma das abordagens que realiza menos simplificações do
transporte de luz, permitindo que diversos efeitos físicos como reflexões e
iluminação indireta surjam naturalmente do algoritmo. Devido ao seu
fotorrealismo, é o algoritmo padrão utilizado para filmes, animações,
simulações e renderizações arquitetônicas. Entretanto, o preço a se pagar por
essa generalidade é o fato de se tratar de um método extremamente custoso
computacionalmente.

Especificamente, \citet{kajiya86}, mostrou um modelo unificado - a equação da
renderização - que modela o transporte de luz numa cena de forma fisicamente
plausível. Algoritmos de Computação Gráfica estimam o valor dessa equação de
modos diversos, simplificando e discretizando partes diferentes da equação. O
\emph{Path Tracing}, por sua vez, produz, utilizando um integrador de Monte
Carlo, estimativas cada vez mais precisas do valor exato da equação. Devido a
essa forma de operação, o tempo gasto pelo algoritmo é aproximadamente
proporcional ao número de estimativas feitas.

Nesse trabalho, exploramos uma técnica, a Amostragem de Importância Múltipla
\citep{veach95}, que permite a criação de estimativas mais próximas do
resultado real, reduzindo a variância do estimador e, consequentemente,
a quantidade de amostras necessárias para uma dada qualidade de imagem. Essa
técnica é a base de algoritmos do estado da arte como ReSTIR
\citep{bitterli20} e VCM \citep{georgiev12}.

Para isso, implementamos um \emph{Path Tracer} e avaliamos o ganho de
qualidade proporcionado pela técnica, comparando imagens renderizadas com igual
número de amostras, com e sem a Amostragem de Importância Múltipla.

Este trabalho está organizado da seguinte forma: o \autoref{cap:fundamentacao}
apresenta os fundamentos matemáticos e físicos necessários, incluindo a equação
da renderização, sua formulação em termos de uma integral de caminhos, o
estimador de Monte Carlo, técnicas de amostragem e a técnica de Amostragem de
Importância Múltipla. O \autoref{cap:implem} descreve a implementação do
\emph{Path Tracer} e a derivação dos pesos da Amostragem de Importância
Múltipla para o caminho completo. Por fim, o \autoref{cap:resultados} apresenta
e analisa os resultados obtidos.


